% concatenated from main.tex, preliminaries.tex, mother_code.tex, additive_channel_proof.tex, conclusion.tex
\documentclass[11pt]{article}
\usepackage{amsmath}
\usepackage{amsfonts}
\usepackage{amssymb}
\usepackage{amsthm}
\usepackage{thmtools, thm-restate}
\usepackage{fullpage}
\usepackage{comment}
\usepackage{hyperref}
\usepackage{lmodern}
\usepackage{caption}
\usepackage{lipsum}
\usepackage{graphicx}
\usepackage{subcaption}
\usepackage{appendix}
\usepackage[normalem]{ulem}


\usepackage[
    backend=biber,
    style=alphabetic,
    sorting=nyt,
    maxbibnames=99 % display all authors
]{biblatex}

\addbibresource{code.bib} %Imports bibliography file

\DeclareMathAlphabet{\mymathbb}{U}{bbold}{m}{n}

\declaretheorem[{style=definition,numberwithin=section}]{definition}
\declaretheorem[{style=definition,sibling=definition}]{theorem}
\declaretheorem[{style=definition,sibling=definition}]{lemma}
\declaretheorem[{style=definition,sibling=definition}]{claim}
\declaretheorem[{style=definition,sibling=definition}]{corollary}
\declaretheorem[{style=definition,sibling=definition}]{proposition}
\declaretheorem[{style=definition,sibling=definition}]{observation}
\declaretheorem[{style=definition,sibling=definition}]{fact}
\declaretheorem[{style=definition,sibling=definition}]{remark}
\declaretheorem[{style=definition,sibling=definition}]{conjecture}
\declaretheorem[{style=definition,name=Question}]{question}


\DeclareMathOperator{\supp}{supp}
\DeclareMathOperator{\dist}{dist}
\DeclareMathOperator{\wt}{wt}
\DeclareMathOperator{\as}{as}
\DeclareMathOperator{\rank}{rank}

\newcommand{\ball}{\mathrm{Ball}}
\newcommand{\bsc}{\mathrm{BSC}}
\newcommand{\weight}{\mathrm{wt}}
\newcommand{\dec}{\mathrm{Dec}}
\newcommand{\enc}{\mathrm{Enc}}

%\newcommand{\basecode}{C_{\text{mo}}}
\newcommand{\basecode}{C_{\text{mother}}}


\newcommand{\typical}{\mathrm{Typical}}

\newcommand{\poly}{\mathrm{poly}}
\newcommand{\polylog}{\mathrm{polylog}}
\newcommand{\quasipoly}{\mathrm{quasi}\textrm{-}\mathrm{poly}}

\newcommand\numberthis{\addtocounter{equation}{1}\tag{\theequation}}

\newtheorem*{theorem*}{Theorem}
\newtheoremstyle{named}{}{}{\itshape}{}{\bfseries}{.}{.5em}{\thmnote{#3}#1}
\theoremstyle{named}
\newtheorem*{namedthm}{}

% Keywords command
\providecommand{\keywords}[1]
{
  % \small	
  \textbf{\textit{Keywords---}} #1
}



\title{
Capacity-Achieving Codes with Inverse-Ackermann-Depth Encoders
% \thanks{
%This work was supported by the CCF-Huawei Populus Grove Fund.
%}
}
\author{
Yuan Li\footnote{Fudan University.  Email: \href{mailto:yuan_li@fudan.edu.cn}{yuan\_li@fudan.edu.cn}.}
}

\date{}

\begin{document}

\maketitle

\begin{abstract}
We prove that for any additive noise channel over $\mathbb{F}_q$, there exist error-correcting codes approaching channel capacity encodable by arithmetic circuits (with weighted addition gates) over $\mathbb{F}_q$ of size $O(n)$ and depth $2\alpha(n)$, where $\alpha(n)$ is a version of the inverse Ackermann function that is at most $3$ for all input lengths $n$ in practice.
Our results demonstrate that certain capacity-achieving codes admit highly efficient encoding circuits that are simultaneously of linear size and inverse-Ackermann depth. 
Our construction composes a linear code with constant rate and relative distance, based on the constructions of G\'{a}l, Hansen, Kouck\'{y}, Pudl\'{a}k, and Viola [IEEE Trans. Inform. Theory 59(10), 2013] and Drucker and Li [COCOON 2023], with an additional layer formed by a disperser graph.
A probabilistic argument over the edge weights of the 
disperser shows the existence of a deterministic encoder achieving error probability $2^{-\Omega(n)}$ at any rate below capacity.
\end{abstract}

\keywords{additive noise channel, error-correcting code, arithmetic circuit, inverse Ackermann function, disperser}

\input{Introduction}

\section{Preliminaries} \label{ch:prelim}

\subsection{Discrete Memoryless Channel}

A \emph{discrete memoryless channel} $\Pi$ consists of input and output alphabets $\mathcal{X}$, $\mathcal{Y}$, and transition probabilities $(p(y \mid x))_{x \in \mathcal{X},\, y \in \mathcal{Y}}$, with $p(y \mid x) = \prod_{i=1}^n p(y_i \mid x_i)$ for $x \in \mathcal{X}^n$, $y \in \mathcal{Y}^n$.

A channel is \emph{symmetric} if $|\mathcal{X}| = |\mathcal{Y}| = q$ and every row of its transition matrix is a permutation of every other row, and likewise for columns. Consequently, every row and every column contains the same multiset of probabilities.

A symmetric channel over $\mathbb{F}_q$ is an \emph{additive noise channel} 
if $\mathcal{X} = \mathcal{Y} = \mathbb{F}_q$ and $p(z \mid x) = p_N(z - x)$ 
for some distribution $p_N$ over $\mathbb{F}_q$, where arithmetic is in 
$\mathbb{F}_q$; equivalently, the channel output is $x + \eta$ where 
$\eta \sim p_N$ is independent of the input.

By Shannon's noisy-channel coding theorem,
$
C(\Pi) = \max_{p(x)} I(X;Y),
$
where $p(x)$ ranges over all distributions on $\mathcal{X}$. For symmetric channels, the maximum is achieved by the uniform input distribution, which induces a uniform output distribution, giving
\begin{equation}
\label{equ:sym_channel_capacity}
C(\Pi) = \log q - H_\Pi,
\end{equation}
where $H_\Pi$ denotes the binary entropy of the first row of the transition matrix.

\subsection{Error-Correcting Code}

Let $C : \mathbb{F}_q^{k} \to \mathbb{F}_q^{n}$. The \emph{rate} of $C$ is $k/n$. Its \emph{minimum distance} is
\[
d(C) = \min_{\substack{x, y \in \mathbb{F}_q^k \\ x \ne y}} \mathrm{dist}(C(x),C(y)),
\]
where $\mathrm{dist}(u,v) = |\{i \in [n] : u_i \ne v_i\}|$ denotes Hamming distance, and its \emph{relative distance} is $d(C)/n$.

Let $\enc:\mathbb{F}_q^{k} \to \mathbb{F}_q^{n}$ be an encoder and $\dec:\mathbb{F}_q^{n} \to \mathbb{F}_q^{k} \cup \{\text{fail}\}$ a decoder. The \emph{failure probability} of $(\enc, \dec)$ over a channel $\Pi$ is
\[
P_e = \max_{m \in \mathbb{F}_q^k} \sum_{y \in \mathbb{F}_q^n} p(y \mid \enc(m))\, \mathbf{1}[\dec(y) \neq m].
\]


\subsection{Arithmetic Circuit}

Fix a finite field $\mathbb{F}_q$. A \emph{linear circuit} over $\mathbb{F}_q$ consists of weighted addition gates of unbounded fan-in, where a gate with inputs $x_1,\dots,x_s$ computes $\sum_{i=1}^s c_i x_i$ for fixed $c_i \in \mathbb{F}_q$; such circuits compute linear functions over $\mathbb{F}_q$. The \emph{size} of a circuit is its total number of \emph{wires}, and its \emph{depth} is the length of its longest input-to-output path.




\subsection{Ackermann Function}

\begin{definition}
\label{def:acker}
(Ackermann function \cite{Tar75, DDP+83})
The \emph{Ackermann function} $A : \mathbb{N} \times \mathbb{N} \to \mathbb{N}$ is defined recursively by
\begin{equation}
A(i, j) \,=\,
\begin{cases}
    2j,                   & i = 0,\ j \ge 1,   \\
    2,                    & i \ge 1,\ j = 1,   \\
    A(i-1,\, A(i, j-1)), & i \ge 1,\ j \ge 2.
\end{cases}
\end{equation}
%We write $A_i(j)$ as shorthand for $A(i,j)$.
\end{definition}

For example, $A(0, j) = 2j$, $A(1, j) = 2^j$ and $A(2, j) = 2\uparrow j$, where $2\uparrow j$ denotes a tower of $j$ twos:
\begin{equation*}
2\uparrow j \;=\; \underbrace{2^{2^{\cdot^{\cdot^{\cdot^{2}}}}}}_{j}.
\end{equation*}

\begin{definition} \cite{RS03}
\label{def:iter}
For a function $f : \mathbb{N} \to \mathbb{N}$, the \emph{iterated function}
$f^* : \mathbb{N} \to \mathbb{N}$ is defined as
\begin{equation}
    f^*(n) \;=\; \min\bigl\{\, k \ge 0 : f^{(k)}(n) \le 1 \,\bigr\},
\end{equation}
where $f^{(0)}(n) = n$ and $f^{(k)}(n) = f\bigl(f^{(k-1)}(n)\bigr)$
for $k \ge 1$.
\end{definition}

\begin{definition}
\label{def:inv-acker} For each integer $d \ge 0$, the function $\lambda_d : \mathbb{N} \to \mathbb{N}$
is defined as
\begin{equation}
    \lambda_d(n) \;=\; \min\bigl\{\, m \ge 1 : A(d,\, m) \ge n \,\bigr\}.
\end{equation}
\end{definition}

For example, \[
\lambda_0(n)=\left\lceil \tfrac{n}{2} \right\rceil,\quad
\lambda_1(n)=\left\lceil \log_2 n \right\rceil,\quad
\lambda_2(n)=\log^* n.
\]


\begin{proposition}
\label{prop:lambda-closed-form}
For all $n \ge 1$,
\begin{align*}
    \lambda_0(n) &\;=\; \left\lceil \frac{n}{2} \right\rceil,  \\
    \lambda_d(n) &\;=\; \lambda^*_{d-1}(n), \qquad d \ge 1.   
\end{align*}
\end{proposition}
\begin{proof}
For $d=0$, since $A(0,m)=2m$, we have
\[
\lambda_0(n)=\min\{m\ge1:2m\ge n\}=\left\lceil \tfrac{n}{2}\right\rceil.
\]
For $d\ge1$, using $A(d,m)=A(d-1,A(d,m-1))$ and the definition of $\lambda_{d-1}$, we have
\[
A(d,m)\ge n
\;\Longleftrightarrow\;
A(d,m-1)\ge \lambda_{d-1}(n).
\]
Iterating this equivalence yields
\[
A(d,m)\ge n
\;\Longleftrightarrow\;
\lambda_{d-1}^{(m-1)}(n)\le 2,
\]
which is equivalent to $\lambda_{d-1}^{(m)}(n)\le 1$. Hence
\[
\lambda_d(n)=\min\{m\ge1:A(d,m)\ge n\}
=\min\{k\ge0:\lambda_{d-1}^{(k)}(n)\le1\}
=\lambda_{d-1}^*(n).
\]
\end{proof}


\begin{definition}
\label{def:alpha}
The \emph{inverse Ackermann function} $\alpha : \mathbb{N} \to \mathbb{N}$ is defined as
\begin{equation}
    \alpha(n) \;=\; \min\bigl\{\, d \ge 0 : \lambda_d(n) \le 3 \,\bigr\}.
\end{equation}
\end{definition}


The function $\alpha(n)$ is well-defined since $A(i, 3) \ge i$ for all $i \ge 1$, so $\lambda_n(n) \le 3$ for any $n \ge 1$, witnessing $d \le n$. The function $\alpha(n)$ grows extremely slowly: for example, $\alpha(n) \le 3$ for all $n \le 2^{65536}$.









\subsection{Dispersers and Superconcentrators}


\begin{definition} \cite{guruswami2012essential} A bipartite graph $G = (L = [n], R = [m], E)$ is a $(\gamma, \varepsilon)$-disperser if for all subsets $S \subseteq L$ with $|S| \ge \gamma n$, we have $|N(S)| \ge (1 - \varepsilon)m$.
\end{definition}

\begin{theorem}[Theorem 1.10 in \cite{RT00}, restated] \label{thm:prob_disperser} Let $c > 0$ and $\gamma, \varepsilon > 0$. For any positive integer $n$ and $m = \lfloor cn \rfloor$, there exists a $(\gamma, \varepsilon)$-disperser graph $G = (V_1 = [n], V_2 = [m], E)$ with degree bounded by $O_{c, \gamma, \varepsilon}(1)$.
\end{theorem}

The original Theorem 1.10 in \cite{RT00} is more general, allowing cases where $m \gg n$. We state a simplified version tailored to our needs. The original theorem guarantees that the left degree is bounded; the right degree is unbounded.
By applying a purging argument that discards the half of the right vertices with the highest degrees, one can also ensure that the right-degree remains bounded.

\begin{definition} \cite{Val75}
A directed acyclic graph on $m$ inputs and $n$ outputs is a \emph{superconcentrator} if for every $k \le \min\{m, n\}$ and every subset $S$ of $k$ inputs and $T$ of $k$ outputs, there exist $k$ vertex-disjoint paths from $S$ to $T$.
\end{definition}



\section{Mother Code Construction}

A linear map $C: \mathbb{F}_q^n \to \mathbb{F}_q^{32n}$ is a \emph{good code} if $\mathrm{wt}(C(x)) \ge 4n$ for every nonzero $x \in \mathbb{F}_q^n$ \cite{drucker2023minimum}, where the constant $32$ follows~\cite{GHK+12}.

A linear map $C : \mathbb{F}_q^{n} \to \mathbb{F}_q^{32n}$ is an \emph{$(n, r, s)$-partial good code}, $(n,r,s)$-PGC for short, if $\mathrm{wt}(C(x)) \ge 4n$ for every nonzero $x \in \mathbb{F}_q^{n}$ with $\mathrm{wt}(x) \in [r, s]$, where $\mathrm{wt}(x) = |\{i : x_i \neq 0\}|$ denotes the Hamming weight of $x$.


\begin{definition} \label{def:range_detector}
\cite{GHK+12} An \emph{$(m, n, \ell, k, r, s)$-range detector} is a map $C:\mathbb{F}_q^m \to \mathbb{F}_q^n$ such that $\mathrm{wt}(C(x)) \in [r, s]$ for every $x \in \mathbb{F}_q^m$ with $\mathrm{wt}(x) \in [\ell, k]$. When $s = n$, we omit the last parameter and write $(m, n, \ell, k, r)$-range detector.
\end{definition}

A good code is a special case of a range detector: an $(n, r, s)$-PGC is an $(n, 32n, r, s, 4n)$-range detector.


Fix a finite field $\mathbb{F}_q$. Let $S_d(n)$ denote the minimum size of a depth-$d$ linear circuit over $\mathbb{F}_q$ computing a good code $C : \mathbb{F}_q^n \to \mathbb{F}_q^{32n}$. More generally, let $S_d(n, r, s)$ denote the minimum size of a depth-$d$ linear circuit over $\mathbb{F}_q$ computing an $(n,r,s)$-PGC $C : \mathbb{F}_q^n \to \mathbb{F}_q^{32n}$.


Our goal is to prove the following theorem. The proof closely follows that of Theorem~1.1 in \cite{drucker2023minimum}, which refines the construction and analysis of \cite{GHK+12}; we therefore present only a proof sketch, highlighting the differences.

\begin{theorem}
\label{thm:mother_code}
Fix a finite field $\mathbb{F}_q$. For any positive integer $n$ and integer $d \ge 0$,
\[
S_{2d}(n) = O_q\!\bigl(\lambda_d(n) \cdot n\bigr).
\]
In particular, taking $d = \alpha(n)$ gives $S_{2\alpha(n)}(n) = O_q(n)$.
\end{theorem}


\subsection{Construction Overview}

We begin by outlining the construction and its main building blocks. In its structure, the entire construction closely resembles superconcentrators \cite{Val75, Val76, Val77, DDP+83}.

Valiant introduced superconcentrators while studying circuit lower bounds. For proving circuit lower bounds, the existence of ultra-low depth, linear-size superconcentrators is a negative result, as it rules out superlinear 
lower bounds based solely on information-transfer arguments. On the other hand, superconcentrators eliminate information-transfer bottlenecks, making them useful in computation and communication tasks, for example, secret sharing \cite{Li23}.

The construction relies on the following key components:
\begin{itemize}
    \item \textbf{Distance amplifier.} A \emph{distance amplifier} can improve the rate and/or the relative distance from any constant to near the Gilbert--Varshamov bound, and it can be computed by a depth-1 circuit of size $O(n)$ with bounded fan-in (Lemma~\ref{lem:rate_booster}).
    
    
    \item \textbf{Output amplifier.} An \emph{output amplifier} arbitrarily increases the number of output coordinates while maintaining a constant relative Hamming weight, realizable by a depth-$1$ circuit of linear size (Lemma~\ref{lem:output_amplifier}).
    
    \item \textbf{Condenser.} A \emph{condenser} reduces the input length 
    from $n$ to $n/r$ while preserving a lower bound on the output Hamming 
    weight, computable by a depth-$1$ circuit of linear size 
    (Lemma~\ref{lem:condenser}).
    
    \item \textbf{Composition Lemma.} Combines several PGCs into a larger one, increasing the depth by 1 while keeping the output fan-in bounded by a constant (Lemma~\ref{lem:composition}).
\end{itemize}

%Using these building blocks, a good code can be constructed recursively with inverse-Ackermann depth and a linear number of wires.



\subsection{Proof of Theorem \ref{thm:mother_code}}

\begin{lemma} [Distance Amplifier]
\label{lem:rate_booster}	
Let $C:\mathbb{F}_q^n \to \mathbb{F}_q^{32n}$ be a code of relative distance $\rho > 0$.
For any $c > 1$ and $\delta > 0$ satisfying $\tfrac{1}{c} < 1 - H_q(\delta)$, there exists a
linear map
\[
L : \mathbb{F}_q^{32n} \to \mathbb{F}_q^{\lfloor cn \rfloor}
\]
such that $L(C(x))$ has relative distance at least $\delta$. Moreover, $L$ is computable by
depth-1 linear circuits of size $O_{\rho,c,\delta}(n)$, and all output gates have bounded
fan-in $O_{\rho,c,\delta}(1)$.
\end{lemma}
\begin{proof}
Identical to that of Lemma~3.4 of \cite{drucker2023minimum}, generalized to $\mathbb{F}_q$.
We take a bounded-degree $(\delta,\varepsilon)$-disperser $G=([32n],[cn],E)$, replace each right vertex in $[cn]$ by a weighted addition gate, and assign to every edge an independently and uniformly chosen coefficient from $\mathbb{F}_q$.

Let $x \in \mathbb{F}_q^n$ be nonzero, and let $N(C(x)) \subseteq [\lfloor cn \rfloor]$ denote the set of output coordinates adjacent in $G$ to at least one nonzero coordinate of $C(x)$.
The distribution of $G(C(x))\restriction_{N(C(x))}$ is uniform in $\mathbb{F}_q^{N(C(x))}$, with randomness arising from the coefficients on the edges. Applying a union bound and using the inequality
\begin{equation}
\label{equ:bsum_fq}
\sum_{i=0}^{\gamma n} \binom{n}{i}(q-1)^i \le q^{H_q(\gamma)n},
\end{equation}
we conclude that there exists such a linear map $L$.
\end{proof}


The following lemma allows us to combine multiple PGCs into a larger one, which will be applied repeatedly in the construction.

\begin{lemma} [Composition Lemma]
\label{lem:composition}
For any $1 \le r_1 < \cdots < r_{t+1} \le n$, we have
\[
S_{d+1}(n, r_1, r_{t+1}) \le \sum_{i=1}^t S_d(n, r_i, r_{i+1}) + O(tn).
\]Furthermore, if, for each $i = 1, \ldots, t$, there exists an $(n, r_i, r_{i+1})$-PGC computable by a depth-$d$ size-$s_i$ linear circuit with output gates of bounded fan-in $D$, then
$
S_d(n, r_1, r_{t+1}) \le \sum_{i=1}^t s_i + O(D t n),
$
and the output gates of the combined circuit have bounded fan-in $O(Dt)$.
\end{lemma}
\begin{proof}
The proof is analogous to Lemma 3.5 in \cite{drucker2023minimum}, but over a general field $\mathbb{F}_q$. We provide a sketch of the argument.

Let $C_i:\mathbb{F}_q^n \to \mathbb{F}_q^{32n}$ denote an $(n, r_i, r_{i+1})$-PGC computable by a linear circuit of size $S_d(n, r_i, r_{i+1})$ and depth $d$. Let $y_1, \ldots, y_{32n}$ be the gates on layer $d+1$, defined by
\[
y_j = \sum_{i=1}^t \alpha_{j,i} \, C_i(x)_j,
\]
where $j = 1, \ldots, 32n$ and the coefficients $\alpha_{j,i} \in \mathbb{F}_q$ are chosen uniformly at random from $\mathbb{F}_q$. A union bound argument shows that there exist coefficients $\alpha_{j,i}$ such that $\mathrm{wt}(y) \ge \frac{n}{4}$ for all $x \in \mathbb{F}_q^n$ with $\weight(x) \in [r_1, r_{t+1}]$.

By applying the rate amplifier, the distance can be increased from $n/4$ to $4n$, producing a circuit of depth $d+2$. Since the rate amplifier's output gates have bounded fan-in $O(1)$, the final layer can be collapsed, resulting in a circuit of depth $d+1$ and size $\sum_{i=1}^t S_d(n, r_i, r_{i+1}) + O(t n).$

Assume that each $(n, r_i, r_{i+1})$-PGC has output gates of bounded fan-in $D$. Collapsing the last layer, we obtain a circuit of depth $d$ and size $\sum_{i=1}^t S_d(n, r_i, r_{i+1}) + O(Dtn)$.
\end{proof}


The following lemma introduces an \emph{output amplifier}, which increases 
the number of output coordinates while preserving a relative distance of 
at least $1/8$.

\begin{lemma} [Output Amplifier] 
\label{lem:output_amplifier}
Fix a finite field $\mathbb{F}_q$. For every positive integer $n$ and every $m \ge 3n$, there exists an 
$(n, m, n/8, n, m/8)$-range detector over $\mathbb{F}_q$ that can be computed by a depth-1 
linear circuit of size $O(m)$. In addition, each output gate has fan-in bounded by an absolute constant.
\end{lemma}
\begin{proof}
Identical to Lemma~3.8 of~\cite{drucker2023minimum}, with the $\mathbb{F}_2$ 
union bound replaced by inequality~\eqref{equ:bsum_fq}.
\end{proof}

%The proof of Lemma \ref{lem:output_amplifier} follows the same argument as Lemma 3.8 of \cite{drucker2023minimum}, using a union bound together with inequality \eqref{equ:bsum_fq}.

\begin{lemma}[Condenser~{\cite{GHK+12}}]
\label{lem:condenser}
Fix a field $\mathbb{F}_q$. There exists a constant $c_0 = c_0(q)$ such 
that for all integers $n$, real numbers $r, s$ with $c_0 \le r \le n$ 
and $s \in [1, n/r^{1.5}]$, there exists an
\[
(n,\, \lfloor n/r \rfloor,\, s,\, n/r^{1.5},\, s,\, \lfloor n/r \rfloor)
\text{-range detector}
\]
computable by a depth-$1$ linear circuit of size $O(n)$.
\end{lemma}
\begin{proof}
The proof closely follows Lemma~23 of~\cite{GHK+12}, with the only 
difference being the verification that for $r \ge c_0(q)$ sufficiently 
large,
\[
\left(\frac{6e\ell}{\frac{5}{6}\cdot m}\right)^{\frac{5}{6}\cdot 6\ell} 
\le 2^{-\ell}(q-1)^{-\ell}\cdot \binom{n}{\ell},
\]
where $m = \lfloor n/r \rfloor$ and $\ell \in [r,\, n/r^{1.5}]$.
\end{proof}

\begin{lemma} [Reduction Lemma]
\label{lem:reduction}
Fix a finite field $\mathbb{F}_q$. Let $c_0 = c_0(q)$ be the constant in Lemma \ref{lem:condenser}.
 For any $r \in [c_0, n]$ and $1 \le s \le t \le \frac{n}{r^{1.5}}$,
\begin{equation}
\label{equ:red_ineq}
S_d(n, s, t) \,\,\le\,\, S_{d-2}\left(\left\lfloor \frac{n}{r} \right\rfloor, s, \frac{n}{r}\right) + O(n).
\end{equation}
In addition, the output gates computing the $(n, s, t)$-PGC have bounded fan-in $O(1)$.
\end{lemma}

The proof of Lemma \ref{lem:reduction} is the same as that of Lemma 3.9 in \cite{drucker2023minimum}.

\begin{lemma} [Depth-$2$ Base Case]
\label{lem:depth2}
Fix $\mathbb{F}_q$. For any $r \in [1, n]$, we have
\[
S_2\left(n, \frac{n}{r}, n\right) = O_q\left(\log^2 r \cdot n\right).
\]
\end{lemma}
\begin{proof}
The proof of Lemma \ref{lem:depth2} follows the same argument as in the proof of Lemma 27 in \cite{GHK+12}; we provide a brief sketch here.

Let $k_1 = r$ and define $k_{i+1} = k_i / 2$. Let $t$ be the smallest integer such that $k_t \le 1$.  
Our strategy is to first construct $(n, n/k_i, n/k_{i+1})$-PGCs of size $O_q(n \log r)$, and then use the Composition Lemma to combine $O(\log r)$ such PGCs.

We construct an $(n, n/k_i, n/k_{i+1})$-PGC as follows. Let $n_i = O\bigl(\frac{n}{k_i} \cdot \log r\bigr)$, and let the middle layer be $y_1, \ldots, y_{n_i}$. Each $y_j$ is connected to $O(k_i)$ inputs chosen independently and uniformly at random with replacement, with coefficients drawn uniformly from $\mathbb{F}_q$. One can argue that $\Pr[y_j = 0] \le 1/2$. Applying a Chernoff bound, we obtain
\[
\Pr\bigl[\weight(y) \le \frac{n_i}{8}\bigr] \le 2^{-\frac{3}{64} \cdot n_i} <
\sum_{j \le n/k_{i+1}} \binom{n}{j}.
\]
A union bound then implies the existence of a linear circuit such that $\weight(y) \ge n_i/8$ for all $y \in \mathbb{F}_q^n$ with $\weight(y) \in [n/k_i, n/k_{i+1}]$. Finally, by placing an output amplifier (Lemma \ref{lem:output_amplifier}) at the bottom, we obtain an $(n, n/k_i, n/k_{i+1})$-PGC of depth $2$, and by Lemma \ref{lem:output_amplifier}, the fan-in of each output gate is bounded by $O(1)$.

Applying the Composition Lemma~\ref{lem:composition} to compose the $t$ 
PGCs, $(n, n/k_i, n/k_{i+1})$ for $i = 1, \ldots, t$, gives an $(n, n/r, n)$-PGC of depth $2$ and size 
$O_q(n \log^2 r)$.
\end{proof}

\begin{lemma} [Depth 4 Base Case]
\label{lem:depth4}
Fix a finite field $\mathbb{F}_q$. For any $r \in [1, n]$, we have
\[
S_4\left(n, \frac{n}{r}, n\right) = O_q\left(\lambda_2(n)\cdot n\right).
\]
\end{lemma}
\begin{proof}
The proof follows Lemma~26 of~\cite{GHK+12}; we sketch the key steps.


Let $c_0 = c_0(q)$ be the constant from Lemma \ref{lem:condenser}.  
If $r < c_0$, then $S_4(n, n/r, n) = O_q(n)$ since a distance amplifier 
(Lemma~\ref{lem:rate_booster}) suffices.
Hence, we may assume $r \ge c_0$.  

Let $k_1 = c_0$ and define $k_{i+1} = 2^{\sqrt{k_i}}$, and let $t$ be the smallest integer such that $k_t \ge n$.  
Note that $k_{i+2} \ge 2^{k_i}$, which implies that $t = O(\log^* n) = O(\lambda_2(n))$.

By Lemma \ref{lem:depth2}, we have
\begin{align*}
S_2\Bigl(\frac{n}{k_i^{2/3}}, \frac{n}{k_{i+1}}, \frac{n}{k_i}\Bigr) 
&= S_2\Bigl(\frac{n}{k_i^{2/3}}, \frac{n}{2^{\sqrt{k_i}}}, \frac{n}{k_i}\Bigr) \\
&\le \frac{n}{k_i^{2/3}} \, O_q\Bigl(\log^2 2^{\sqrt{k_i}}\Bigr) \\
&= O_q(n).
\end{align*}
Applying Lemma \ref{lem:reduction}, we then obtain an $(n, n/k_{i+1}, n/k_i)$-PGC computable by a depth-$4$ linear circuit with output fan-in bounded by $O(1)$.

Finally, by applying the Composition Lemma (Lemma \ref{lem:composition}) to combine the $O(\log^* n)$ PGCs, we obtain a $(n, n/r, n)$-PGC of depth $4$ and size $O_q(n \log^* n) = O_q(\lambda_2(n) \cdot n)$.
\end{proof}


The following theorem provides the main construction and is proved by induction on $k$. It is established for $q = 2$ in Theorem~3.13 of~\cite{drucker2023minimum}. The argument extends to a general finite field $\mathbb{F}_q$; we omit the details.

\begin{theorem}
\label{thm:ub_main}
Fix a finite field $\mathbb{F}_q$. Let $c_0 = c_0(q)$ be the constant 
from Lemma~\ref{lem:condenser}. There exist constants $c, D > 0$, 
depending only on $q$, such that the following statements hold.
\begin{enumerate}
    \item For any $c_0 \le r \le n$ and any $k \ge 3$,
    \begin{equation}
    \label{equ:d2k_first}
    S_{2k}\!\left(n,\, \frac{n}{A(k-1,r)},\, \frac{n}{r}\right) \le 2cn,
    \end{equation}
    and the output gates of the corresponding linear circuit have fan-in 
    bounded by $D$.
    \item For any $2 \le r \le n$ and any $k \ge 2$,
    \begin{equation}
    \label{equ:d2k_second}
    S_{2k}\!\left(n,\, \frac{n}{r},\, n\right) \le 3c\,\lambda_k(r) \cdot n.
    \end{equation}
    The linear circuits encoding the $(n, n/r, n)$-PGC do not necessarily 
    have bounded output fan-in.
\end{enumerate}
\end{theorem}

Theorem \ref{thm:ub_main} immediately implies Theorem \ref{thm:mother_code}.



\section{Existence of Capacity-Achieving Codes}


\subsection{Code Construction}
\label{section:code_construction}

Let $H = (L=[32k], R = [n], E)$ be a $(1/8, \gamma)$-disperser graph for some small constant $\gamma > 0$. Recall that the disperser property guarantees $|N(S)| \ge (1-\gamma)n$ for every subset $S \subseteq L$ with $|S| \ge |L|/8$. Given an edge-labeling $\alpha: E(H) \to \mathbb{F}_q$, define the linear map $D_{H,\alpha}: \mathbb{F}_q^{32k} \to \mathbb{F}_q^n$ by
\begin{equation}
    D_{H, \alpha}(x_1, \dots, x_{32k})_j = \sum_{(i,j) \in E(H)} \alpha(i,j)\, x_i,
\end{equation}
where all arithmetic is over $\mathbb{F}_q$.


Let $\basecode:\mathbb{F}_q^k \to \mathbb{F}_q^{32k}$ be a linear code with minimum distance at least $4k$. 
The encoder $\enc:\mathbb{F}_q^k \to \mathbb{F}_q^n$ is defined by
\begin{equation}
    \enc(x) = D_{H, \alpha}\bigl(\basecode(x)\bigr).
\end{equation}

Let $\typical_\epsilon(y)$ denote the set of typical transmitted vectors corresponding to the received word $y \in \mathbb{F}_q^n$ for a small constant $\epsilon > 0$  (see Definition~\ref{def:typical_set} below). The decoder outputs the unique codeword in $\typical_\epsilon(y)$ if it exists, and \textsf{fail} otherwise. Since our goal is an existential result, decoding efficiency is not a concern.



\subsection{Typical Set}


For a received vector $y$, its conditional typical set consists of those vectors $x$ for which the channel behavior --- transmitting $x$ and receiving $y$ --- is \emph{statistically typical}, i.e., consistent with what one would normally expect from the channel.

\begin{definition}[Conditional Typical Set]
\label{def:typical_set}
Let $\Pi$ be a strongly symmetric channel. For a received sequence $y \in \mathbb{F}_q^n$, define the \emph{conditional typical set} as
\[
\mathrm{Typical}_\epsilon(y) = \left\{ x \in \mathbb{F}_q^n :
p(y \mid x) > 0
\quad \text{and} \quad
\left| -\frac{1}{n} \sum_{i=1}^n \log p(y_i \mid x_i) - H_\Pi \right| \leq \epsilon
\right\},
\]
where $H_\Pi$ is the per-row output entropy of $\Pi$.
\end{definition}

\begin{proposition}
\label{prop:typical_volume}
For any $y \in \mathbb{F}_q^n$ and any $\epsilon > 0$,
\[
|\mathrm{Typical}_\epsilon(y)| \leq 2^{n(H_\Pi + \epsilon)}.
\]
\end{proposition}

\begin{proof}
By definition of $\mathrm{Typical}_\epsilon(y)$, every $x \in \mathrm{Typical}_\epsilon(y)$ satisfies
\[
p(y \mid x) \geq 2^{-n(H_\Pi + \epsilon)}.
\]
Therefore,
\[
1 = \sum_{x \in \mathbb{F}_q^n} p(y \mid x) \geq \sum_{x \in \mathrm{Typical}_\epsilon(y)} p(y \mid x) \geq |\mathrm{Typical}_\epsilon(y)| \cdot 2^{-n(H_\Pi + \epsilon)},
\]
which gives $|\mathrm{Typical}_\epsilon(y)| \leq 2^{n(H_\Pi + \epsilon)}$.
\end{proof}


\begin{lemma}[Chernoff bound]
Let $X_1, X_2, \ldots, X_n \in \{0,1\}$ be independent random variables with
$\mathbb{E}[X_i] = p$, and let $S_n = \sum_{i=1}^n X_i$. 
Then for any $\epsilon > 0$,
\[
\Pr\bigl[\, |S_n - pn| \ge \epsilon n \,\bigr]
\le 2 \exp\!\left( - \frac{\epsilon^2 pn}{2} \right).
\]
\end{lemma}

The following lemma is a quantitative version of the Asymptotic Equipartition Property \cite{CoverThomas2006}; we include the proof for completeness as the exponential bound will be needed later.

\begin{lemma}
\label{lem:typical_prob}
For any $x \in \mathbb{F}_q^n$ and $\epsilon > 0$,
\[
\Pr_{y \sim p(\cdot \mid x)}\!\bigl[x \notin \typical_\epsilon(y)\bigr] \;\le\; 2q\cdot 2^{-\Omega_\Pi(\epsilon^2 n)}.
\]
\end{lemma}
\begin{proof}
Since $\Pi$ is strongly symmetric, the distribution of $Z_i = -\log p(y_i \mid x_i)$
is the same for any fixed $x_i$, and the entropy is the same for every
row. Therefore, without loss of generality, assume $x = 0^n$, so that
$y_1, \ldots, y_n$ are i.i.d.\ with common distribution $p(\cdot \mid 0)$.

Write $p_a = p(a \mid 0)$ for $a \in \mathbb{F}_q$, and for each $a$ let
$X_i^{(a)} = \mathbf{1}[y_i = a]$, so that
\[
\frac{1}{n}\sum_{i=1}^n Z_i
= \sum_{a \in \mathbb{F}_q} (-\log p_a)\cdot \frac{1}{n}\sum_{i=1}^n X_i^{(a)}.
\]
Taking expectations gives $\mathbb{E}[Z_i] = H_\Pi$, so
\[
\frac{1}{n}\sum_{i=1}^n Z_i - H_\Pi
= \sum_{a \in \mathbb{F}_q}(-\log p_a)
  \left(\frac{1}{n}\sum_{i=1}^n X_i^{(a)} - p_a\right).
\]
If the left-hand side exceeds $\epsilon$ in absolute value, then by the triangle
inequality there exists some $b \in \mathbb{F}_q$ with
\[
\left|\frac{1}{n}\sum_{i=1}^n X_i^{(b)} - p_{b}\right|
> \frac{\epsilon}{q \cdot \max_b |\log p_b|} =: \epsilon' > 0.
\]
Since $\{X_i^{(a)}\}_i$ are i.i.d.\ Bernoulli$(p_a)$, the Chernoff bound and a
union bound over all $q$ symbols give
\[
\Pr\!\left[\,\left|\frac{1}{n}\sum_{i=1}^n Z_i - H_\Pi\right| > \epsilon\,\right]
\leq \sum_{a \in \mathbb{F}_q} 2\exp\!\left(-\frac{(\epsilon')^2 p_a\, n}{2}\right)
\leq 2q \cdot \exp(-c_\Pi\,\epsilon^2 n),
\]
where 
\[
c_\Pi = \frac{\min_a p_a}{2(q \cdot \max_a |\log p_a|)^2} > 0
\]
depends only on $\Pi$. Therefore,
\[
\Pr_{y \sim p(\cdot \mid x)}\!\bigl[x \notin \mathrm{Typical}_\epsilon(y)\bigr]
\leq 2q \cdot 2^{-\Omega_\Pi(\epsilon^2 n)}.
\qedhere
\]
\end{proof}



\subsection{Bounding the Error Probability}

In this section, we prove Theorem~\ref{thm:main}. The encoder composes two layers: a \emph{mother code} $\basecode: \mathbb{F}_q^k \to \mathbb{F}_q^{32k}$ with minimum distance $4k$, and a \emph{disperser layer} $D: \mathbb{F}_q^{32k} \to \mathbb{F}_q^n$. The disperser layer is a bipartite disperser graph with edge weights drawn uniformly and independently at random from $\mathbb{F}_q$.
For any nonzero $m \in \mathbb{F}_q^k$, the mother code guarantees $|\supp(\basecode(m))| \ge 4k$; the disperser layer then ensures that $D(\basecode(m))$ is uniformly distributed over a subspace of $\mathbb{F}_q^n$ of dimension at least $(1-\gamma)n$, for a small constant $\gamma > 0$. This mimics the behavior of a random linear code, which is known to achieve channel capacity.

The additive noise structure of $\Pi$ is used in two places. First, it 
ensures that the typical set has the correct volume and probability 
(Proposition~\ref{prop:typical_volume} and Lemma~\ref{lem:typical_prob}). 
Second, it implies that $P_e(m) = P_e(0)$ for every message $m \in 
\mathbb{F}_q^k$ (Claim~\ref{claim:symmetry}), so that it suffices to 
bound the error probability for a single message.

\begin{proof}[Proof of Theorem \ref{thm:main}]
Let $\basecode:\mathbb{F}_q^k \to \mathbb{F}_q^{32k}$ be a fixed linear code 
with minimum distance at least $4k$ (Theorem~\ref{thm:mother_code}). Let $H = (L=[32k],\, R=[n],\, E)$ be a fixed $(1/8,\gamma)$-disperser graph, where $\gamma > 0$ is a small constant to be determined later.

Let $\alpha:E(H) \to \mathbb{F}_q$ assign a field element to each edge. Define the linear map $D_{H,\alpha}: \mathbb{F}_q^{32k} \to 
\mathbb{F}_q^n$ by
\[
    D_{H,\alpha}(x)_j \;=\; \sum_{(i,j)\,\in\, E(H)} \alpha(i,j)\, x_i.
\]

The encoder $\enc:\mathbb{F}_q^k \to \mathbb{F}_q^n$ is defined as the 
two-step composition
\[
    \enc(x) \;=\; D_{H,\alpha}\!\bigl(\basecode(x)\bigr),
\]
where the message is first expanded to $\mathbb{F}_q^{32k}$ via $\basecode$, then mapped 
to $\mathbb{F}_q^n$ via $D_{H,\alpha}$.

\textbf{Proof strategy.} Rather than constructing a good code explicitly, we establish its 
\emph{existence} via a probabilistic argument. Specifically, we fix the  bipartite graph $H$ and randomize only the edge labeling: let $\dot\alpha: E(H) \to \mathbb{F}_q$ be chosen 
uniformly at random, inducing a random encoder
\[
    \dot\enc(x) \;=\; D_{H,\dot\alpha}\!\bigl(\basecode(x)\bigr).
\]
This is equivalently described by a random generator matrix $\dot G = D_{H,\dot\alpha}\,\basecode$, so that $\dot\enc(x) = \dot G x.$

It then suffices to show that the \emph{expected} decoding error under this 
random choice of $\dot G$ is small. By a standard averaging argument, this implies the existence of at least one fixed choice of $\alpha$ --- and hence a deterministic encoder $\enc$ --- that achieves the desired error guarantee.

Let $\dot G \in \mathbb{F}_q^{n \times k}$ be the random generator matrix. Let $m \in \mathbb{F}_q^k$ be the message, and $c = \dot G m \in \mathbb{F}_q^n$ be the codeword.

\begin{claim}
\label{claim:symmetry}
Let $\Pi$ be an additive noise channel over $\mathbb{F}_q$ and let 
$G \in \mathbb{F}_q^{n \times k}$ be any fixed generator matrix. Then 
$P_e(m) = P_e(0)$ for every $m \in \mathbb{F}_q^k$.
\end{claim}
\begin{proof} (of Claim \ref{claim:symmetry})
Define the bijection $\sigma : \mathbb{F}_q^n \to \mathbb{F}_q^n$ by 
$\sigma(z) = z - Gm$, and write $z' = 
\sigma(z)$. Since $\sigma$ is a bijection, substituting 
$z \mapsto z'$ in a sum over $\mathbb{F}_q^n$ is valid. 
We verify three facts.

\medskip
\noindent\textbf{Fact 1.} $p(z \mid Gm) = p(z' \mid 0)$.

Since $\Pi$ is an additive noise channel, $p(z_i \mid x_i) = p_N(z_i - x_i)$, 
hence by independence across coordinates,
\[
p(z \mid Gm) 
= \prod_{i=1}^n p_N(z_i - (Gm)_i) 
= \prod_{i=1}^n p_N(z_i') 
= p(z' \mid 0).
\]

\medskip
\noindent\textbf{Fact 2.} $\mathrm{Typical}_\epsilon(z) - Gm = 
\mathrm{Typical}_\epsilon(z')$.

For any $x \in \mathbb{F}_q^n$, by the additive noise structure,
\[
p(z_i \mid x_i) = p_N(z_i - x_i) = p_N(z_i' - (x_i - (Gm)_i)) 
= p(z_i' \mid x_i - (Gm)_i),
\]
hence
\[
-\frac{1}{n}\sum_{i=1}^n \log p(z_i \mid x_i) 
= -\frac{1}{n}\sum_{i=1}^n \log p(z_i' \mid x_i - (Gm)_i).
\]
Therefore $x \in \mathrm{Typical}_\epsilon(z) \iff 
x - Gm \in \mathrm{Typical}_\epsilon(z')$, which gives 
$\mathrm{Typical}_\epsilon(z) - Gm = \mathrm{Typical}_\epsilon(z')$. In particular $Gm \in \mathrm{Typical}_\epsilon(z) \iff 
0 \in \mathrm{Typical}_\epsilon(z')$.

\medskip
\noindent\textbf{Fact 3.} $\exists\, m' \neq m : Gm' \in 
\mathrm{Typical}_\epsilon(z) \iff \exists\, \hat{m} \neq 0 : 
G\hat{m} \in \mathrm{Typical}_\epsilon(z')$.

By Fact 2, $Gm' \in \mathrm{Typical}_\epsilon(z) \iff Gm' - Gm 
\in \mathrm{Typical}_\epsilon(z')$. Setting $\hat{m} = m' - m$ 
and using linearity of $G$, we have $Gm' - Gm = G\hat{m}$, and 
$m' \neq m \iff \hat{m} \neq 0$.

\medskip
Therefore, applying the substitution $z \mapsto 
z' = z - Gm$ and the three facts above, we have
\begin{align*}
P_e(m) 
&= \sum_{z} p(z \mid Gm)\cdot\mathbf{1}\!\Big[
    Gm \notin \mathrm{Typical}_\epsilon(z)\ \vee\ 
    \exists\, m' \neq m : Gm' \in \mathrm{Typical}_\epsilon(z)
\Big] \\
&= \sum_{z'} p(z' \mid 0)\cdot\mathbf{1}\!\Big[
    0 \notin \mathrm{Typical}_\epsilon(z')\ \vee\ 
    \exists\, \hat{m} \neq 0 : G\hat{m} \in \mathrm{Typical}_\epsilon(z')
\Big] \\
&= P_e(0). \qedhere
\end{align*}
\end{proof}

Claim~\ref{claim:symmetry} reduces the analysis to bounding $P_e(0)$; it thus suffices to exhibit $G$ for which $P_e(0)$ is small.

\noindent\textbf{Step 1: Expressing the error probability.}  
By the decoder defined in Section~\ref{section:code_construction}, the error event for message $0 \in \mathbb{F}_q^k$ consists of either a typicality failure or a collision with another codeword:
\[
P_e(0) = \sum_{z \in \mathbb{F}_q^n} p(z \mid 0) \, 
\mathbf{1}\Big[ 0 \notin \typical_\epsilon(z) \ \vee \ \exists m \neq 0 : \dot G m \in \typical_\epsilon(z) \Big].
\]
By linearity of expectation over a random generator matrix $\dot G = \basecode C_{H, \dot\alpha}$ (with $\alpha : E(H) \to \mathbb{F}_q$ uniform), we have
\[
\mathbb{E}_{\dot G}[P_e(0)]
\le \sum_{z} p(z \mid 0) \mathbf{1}[0 \notin \typical_\epsilon(z)] 
+ \sum_{m \neq 0} \sum_{z} p(z \mid 0) \Pr[\dot G m \in \typical_\epsilon(z)].
\]
Define
\[
E_1 = \sum_{z} p(z \mid 0) \mathbf{1}[0 \notin \typical_\epsilon(z)], \qquad
E_2 = \sum_{m \neq 0} \sum_{z} p(z \mid 0) \Pr[\dot G m \in \typical_\epsilon(z)].
\]

\noindent\textbf{Step 2: Bounding $E_1$.}  
By Lemma \ref{lem:typical_prob}, we have $E_1 \le 2q \cdot 2^{-\Omega_\Pi(\epsilon^2 n)}$.

\noindent\textbf{Step 3: Bounding $E_2$.}  
We bound the probability that some nonzero $m \in \mathbb{F}_q^k$ produces a codeword in $\typical_\epsilon(z)$.

Since $\basecode$ has minimum distance at least $4k$, every nonzero $m$ satisfies
$|\supp(\basecode(m))| \ge 4k = |L|/8$.
By the
$(1/8,\gamma)$-disperser property of $H$, there exists $S \subseteq [n]$ with
$|S| \ge (1-\gamma)n$ such that every $v \in S$ has at least one neighbor
$u \in \supp(\basecode(m))$. For each such $v$, fix any such edge $e=(u,v)$:
since $\basecode(m)_u \neq 0$ and $\dot\alpha(e)$ is uniform over $\mathbb{F}_q$,
the output $(\dot G m)_v$ is uniform over $\mathbb{F}_q$. Independence of edge
weights across edges implies these outputs are independent across $v \in S$, so
\[
    (\dot G m)\restriction_S \;\sim\; \mathrm{Uniform}(\mathbb{F}_q^S).
\]

By Proposition~\ref{prop:typical_volume}, $|\typical_\epsilon(z)| \le 2^{n(H_\Pi+\epsilon)}$, hence $|\typical_\epsilon(z)\restriction_S| \le 2^{n(H_\Pi+\epsilon)}$. Thus,
\begin{align*}
    \Pr\bigl[\dot G m \in \typical_\epsilon(z)\bigr]
    &\le \Pr\bigl[(\dot G m)\restriction_S \in \typical_\epsilon(z)\restriction_S\bigr] \\
    &\le \frac{2^{n(H_\Pi+\epsilon)}}{q^{(1-\gamma)n}}.
\end{align*}
Summing over all $q^k - 1 \le q^k$ nonzero $m$ gives
\[
    E_2 \;\le\; \frac{q^k \cdot 2^{n(H_\Pi+\epsilon)}}{q^{(1-\gamma)n}}
    \;=\; 2^{\,n\left(r + H_\Pi + \epsilon - (1-\gamma)\log q\right)},
\]
where $r = \frac{k}{n}\log q$ is the code rate. The exponent equals
\begin{equation}
\label{equ:error_exponent}
    n\bigl(r - \log q + H_\Pi + \gamma\log q + \epsilon\bigr).
\end{equation}
For any fixed rate $r < \log q - H_\Pi$ (i.e., below the capacity
\eqref{equ:sym_channel_capacity}), choosing $\epsilon, \gamma > 0$ sufficiently small
makes the exponent $-\Omega_{\Pi,r}(n)$.

\noindent\textbf{Step 4: Conclusion.}  
Combining the bounds for $E_1$ and $E_2$, we conclude that
\[
\mathbb{E}_{\dot G}[P_e(0)] \le E_1 + E_2 \le 2^{- \Omega_{\Pi, r}(n)}.
\]
By averaging, there exists $G$ such that $P_e(0) \le 2^{-\Omega_{\Pi,r}(n)}$. 
Claim~\ref{claim:symmetry} then gives $P_e(m) = P_e(0)$ for all $m \in \mathbb{F}_q^k$. Thus, the worst-case error probability is bounded by $2^{-\Omega_{\Pi,r}(n)}$.

\textbf{Circuit size and depth.} By Theorem~\ref{thm:mother_code}, the mother code $\basecode$ is encodable by a linear circuit of depth $2\alpha(k)$ and size $O_q(k)$, where $k = rn$.
The disperser code $D_{H,\alpha}$ is computed by a linear circuit of depth $1$ and size $O(n)$ with bounded output degree $O_\gamma(1) = O_{\Pi,r}(1)$ (Theorem~\ref{thm:prob_disperser}). Composing these and collapsing the final layer yields a linear circuit of depth $2\alpha(n)$ and size $O_{\Pi,r}(n)$.

\end{proof}



\section{Conclusion}

We have shown that for additive noise channels over $\mathbb{F}_q$ with $q$ a prime power, there exist capacity-achieving codes encodable by linear circuits over $\mathbb{F}_q$ of linear size and inverse-Ackermann depth. 
While the noisy coding theorem demonstrates the abundance of capacity-achieving codes, our result shows that some of them admit encoders of extremely low circuit complexity.

This leaves a natural open problem: to construct error-correcting codes that achieve capacity, admit linear-size circuits of inverse-Ackermann depth for encoding, and also support efficient decoding.

\subsection*{Acknowledgements}

This work was supported by the CCF-Huawei Populus Grove Fund.




\printbibliography

\appendix

\end{document}
